% =============================================================
% Chapitre 1 — Contexte, entreprise et problématique (~5.5 pages)
% =============================================================
\chapter{Contexte, entreprise et problématique}
\label{ch:contexte}

\section{Présentation de MENAL et du service d'accueil}

MENAL-SARL est une société de droit mauritanien, également implantée en France, fondée en
2025 et spécialisée dans les services informatiques, les technologies numériques et la
fourniture de solutions IT. L'entreprise accompagne les acteurs publics et privés dans leur
transformation numérique, en particulier sur les enjeux de digitalisation, de sécurité des
systèmes d'information et de gestion des infrastructures, avec l'ambition de devenir un
acteur de référence en Mauritanie dans le domaine des technologies de l'information.

Son activité s'articule autour de deux axes complémentaires. D'une part, des services
informatiques couvrant l'ensemble du cycle de vie des systèmes d'information : conception
et déploiement d'architectures cloud (Google Cloud Platform, AWS), mise en place de
solutions sécurisées selon les approches DevSecOps et Zero Trust, développement
d'applications web et d'API sécurisées, conception et industrialisation de pipelines de
données, intégration de solutions de supervision et de monitoring, ainsi qu'audit,
sécurisation et optimisation de systèmes existants. D'autre part, la fourniture de
matériels et de solutions informatiques clés en main : équipements (serveurs, postes de
travail, matériel réseau), installation, déploiement d'infrastructures et maintenance.
MENAL-SARL se positionne ainsi comme un prestataire global, assurant un accompagnement de
bout en bout, de la conception à l'exploitation.

\begin{keybox}[title=Lien avec le PFE]
Le présent projet s'inscrit directement dans cet axe stratégique : il répond à la volonté
de MENAL de développer des solutions cloud sécurisées, à travers la conception et le
déploiement d'une architecture Zero Trust sur Google Cloud Platform, intégrant des
pipelines de données, une API sécurisée et une application web, selon une approche
DevSecOps.
\end{keybox}

\begin{attentedoc}
Présentation du \emph{service d'accueil} encore manquante : équipe, mission du service,
positionnement dans l'organigramme de MENAL. Le guide pédagogique demande explicitement de
\emph{« se focaliser sur une présentation personnelle du service dans lequel a été effectué
le projet »} --- ce volet ne peut pas être rédigé à partir du seul document de présentation
générale de l'entreprise déjà fourni.
\end{attentedoc}

\section{Contexte du stage et mission confiée}

Le stage s'inscrit dans une démarche de MENAL visant à héberger plusieurs applications
métier sur une infrastructure Google Cloud Platform (GCP) unique, mutualisée et sécurisée.
L'application \textbf{ELSON}, plateforme éducative de crowdsourcing, sert d'application
exemple pour valider concrètement le socle : c'est un cas réel plutôt qu'un scénario de
laboratoire abstrait.

\begin{keybox}[title=Périmètre du projet]
Le périmètre du PFE porte sur le \textbf{socle d'hébergement} --- point d'entrée, identité,
exécution des services, réseau, données, enrichissement par IA, supervision --- et sur la
\textbf{chaîne de livraison DevSecOps} qui l'alimente. Le code applicatif interne d'ELSON
est explicitement hors périmètre : seule sa capacité à être hébergée par le socle est
étudiée.
\end{keybox}

\todo{Décrire précisément le rôle tenu durant le stage : conception, implémentation, tests
--- avec des exemples concrets de tâches réalisées personnellement (cf. guide pédagogique,
exigence de \og description des tâches quotidiennes \fg{}).}

\section{Problématique}
\label{sec:problematique}

Deux constats motivent ce travail. D'une part, un modèle de sécurité fondé sur la confiance
accordée au réseau interne devient difficile à justifier dès lors que les applications sont
hébergées dans le cloud et que la chaîne de livraison logicielle constitue elle-même une
cible. D'autre part, la conformité à un modèle de sécurité (Zero Trust, DevSecOps) est
souvent \emph{déclarée} sur la base d'un schéma d'architecture, sans qu'elle soit
effectivement démontrée par des tests.

La problématique retenue pour ce PFE est formulée comme suit :

\begin{quote}
\itshape
Comment concevoir et déployer, avec les moyens et le calendrier d'un projet de fin
d'études, un socle d'hébergement conforme aux principes Zero Trust (NIST SP~800-207) où
chaque contrôle de sécurité peut être \textbf{prouvé par un test de contournement qui
échoue} plutôt qu'affirmé sur plan, tout en conservant un projet réaliste --- simplicité
assumée, coûts maîtrisés, écarts entre conception et implémentation documentés plutôt que
masqués ?
\end{quote}

\begin{attentedoc}
Formulation à confirmer avec l'encadrant académique et l'encadrant entreprise avant de la
figer --- elle conditionne la cohérence de l'introduction générale et de la conclusion.
\end{attentedoc}

\section{Étude de l'existant}
\label{sec:etude-existant}

Contrairement à un projet décisionnel qui remplace un système d'information existant, ce
projet ne remplace pas un logiciel : il construit un socle là où aucun n'existait.
\emph{« L'existant »} désigne donc ici l'état des pratiques de sécurité et d'hébergement
avant le socle, tel que documenté par l'audit de sécurité préalable de l'application ELSON.

\subsection{Description de l'existant}

Avant le projet, chaque application est déployée et sécurisée au cas par cas, sans garantie
de cohérence entre applications. ELSON, première application destinée à être hébergée sur le
socle, a fait l'objet d'un audit de sécurité préalable qui a mis en évidence une dette
technique caractéristique d'un environnement construit sans socle commun. Aucun mécanisme de
détection ou de corrélation d'événements de sécurité n'existe avant ce projet, et
l'infrastructure existante n'est pas décrite en code.

\subsection{Critique de l'existant}

Cinq manques structurent directement les objectifs retenus pour ce PFE (§\ref{sec:objectifs}) :

\begin{itemize}
  \item des secrets mal gérés, sans gestion centralisée ;
  \item l'absence de séparation d'identité par service, rendant possible une élévation de
        privilège implicite ;
  \item l'absence de détection --- aucune visibilité sur les menaces avant le projet ;
  \item l'absence de reproductibilité --- une infrastructure non versionnée, donc non
        reconstructible de façon fiable en cas d'incident ;
  \item l'absence de modèle multi-tenant --- impossible d'accueillir une deuxième application
        sans dupliquer tout l'effort de sécurisation.
\end{itemize}

Chacun de ces constats motive directement l'un des cinq premiers objectifs du projet, détaillés
en §\ref{sec:objectifs} (cf. cahier des charges §2.2, objectifs O1 à O5 --- le sixième
objectif, O6, relève de la valeur pédagogique de la soutenance et non d'un manque technique de
l'existant).

\section{Objectifs du PFE}
\label{sec:objectifs}

\begin{enumerate}
  \item Concevoir une architecture cible conforme aux principes Zero Trust (NIST
        SP~800-207) et documenter chaque décision d'architecture dans un registre dédié.
  \item Mettre en place une chaîne de livraison DevSecOps avec des portes de sécurité
        bloquantes (secrets, analyse statique, vulnérabilités des images).
  \item Déployer un système de détection (SIEM) enrichi par un modèle sémantique non
        génératif reliant les vulnérabilités détectées dans le pipeline aux techniques
        d'attaque effectivement observées.
  \item Valider l'architecture par un plan de tests de contournement démontrant chaque
        contrôle de sécurité annoncé.
  \item Documenter, sans les masquer, les écarts entre l'architecture visée et
        l'architecture réellement déployée.
\end{enumerate}

\section{Démarche et plan du rapport}

\todo{Paragraphe de transition annonçant les chapitres 2 à 5 --- à rédiger une fois la
structure définitive validée (cf. \texttt{00\_SOMMAIRE\_PFE\_MENAL.md}).}
